% SOURCE docs/release/2026-10-29/source/docs/writing-repair/2026-10-29/src/1de93e537048440f.tex CHAR 0
\chapter{A narrower policy interval can be more confidently wrong}
\label{gap29:policy-uncertainty}
The servicing example leaves a natural question. If a fitted transition
model uses the observed histories more efficiently than weighting, can it
give an informative interval for the gain from intervention? The answer
depends on why outcomes are missing and whether the state equation is
adequate. Resampling quantifies variation in an estimated model; it does
not establish the model's economic assumptions.

\section{Follow the customer through the calculation}
The follow-up study uses the same three-decision economy as
Chapter~\ref{gap29:servicing}. It selects a threshold policy on 2,000
training customers and evaluates that fixed choice on 1,000 independent
customers. Each of 199 bootstrap draws resamples whole evaluation
customers, including their observation histories, and refits both
transition coefficients and innovation variance.

For a plug-in value estimate $\widehat V$ and bootstrap quantiles
$q^*_{.025},q^*_{.975}$, the basic interval is
$[2\widehat V-q^*_{.975},\,2\widehat V-q^*_{.025}]$.
The improvement over never acting is calculated within the same simulated
paths and resampled customers before forming its interval. Pairing retains
the covariance between policy values. No separate censoring model is fitted:
under the baseline's conditional independence, the observed transition
regression can estimate the common transition law. The adverse dropout case
violates precisely that premise.

The selected policy is held fixed during evaluation resampling. Its target
is the value of the policy selected on that separate training sample.
Repeating selection across independent experimental replications changes
the target accordingly, and each result is compared with the reference for
its selected policy. This is different from estimating the performance of
an entire future retraining procedure.

The study uses 100 independent replications in each of three cases, with a
separate million-trajectory reference. The initial-state distribution and
reward equation are supplied by the construction. Their estimation
uncertainty is consequently absent. Numerical integration uses 8,192
scrambled low-discrepancy normal paths shared across fits and bootstrap
draws. Eight separate scrambles at the true parameters check sensitivity
to numerical integration; parameter uncertainty remains a separate source
of variation. Across those eight scrambles, standard deviations are below
0.010 reward units for values and 0.009 for improvements. This observed
sensitivity is small relative to the statistical widths reported here; it
is not a uniform integration-error bound for every fitted parameter.

\section{Precision under a maintained assumption}
For the separately selected policy, the baseline interval for improvement
has empirical coverage 94\%, with Monte Carlo standard error 2.4 percentage
points. Its average width is 0.711 reward units. The matched weighting
interval has coverage 98\% and width 6.905. The difference in width is
descriptive evidence of what the parametric restrictions buy; the sample
of replications does not establish exact calibration or general superiority.
Three bootstrap coverage decisions lie near the separate reference's
numerical uncertainty band.

A positive lower bound for improvement occurs in 79\% of baseline fitted
analyses and 2\% of the weighted analyses. This is a decision frequency,
not by itself a measure of safety. The selected list includes an unprofitable
threshold in seven of the 100 training samples. Evaluation must preserve
the selected target and the possibility of a poor selection.

\section{Why the apparently useful interval can fail}
\begin{figure}[htbp]
\centering
\includegraphics[width=.95\textwidth]{scientific-gap-closure/2026-10-29/policy-coverage.pdf}
\caption{Coverage of the separately selected policy's improvement, 100
independent replications per case. Bars are Wilson intervals for Monte Carlo
coverage uncertainty. Numerical-reference uncertainty is reported separately.
Narrow fitted intervals rely on correct observation and transition assumptions.}
\end{figure}
When dropout depends on the unseen next shock, the fitted value interval
covers its reference in none of the 100 replications. The fitted interval
for improvement covers in only 6\%, with Monte Carlo standard error
2.4 percentage points. Its bias is $+1.175$ (Monte Carlo standard error
$0.052$), and its average width is 1.140. Every fitted analysis has a
positive lower bound for improvement, including selections of an
unprofitable threshold. A compelling-looking decision can therefore be
an artefact of the observation process.

The weighted improvement interval covers in 96\% in that case, but has
bias about $+0.498$ and average width 6.348. Its value interval covers only
73\%. Cancellation and imprecision in a difference can mask biased value
estimates. The apparently satisfactory improvement coverage does not
validate the false observation model.

Omitting $X$ from the transition regression produces a different failure:
coverage for improvement falls to 61\%, with Monte Carlo standard error
4.9 percentage points. The estimated bias is $-0.253$ (Monte Carlo standard
error $0.019$), despite an average interval width of 0.689.
Resampling the fitted omission repeats it.

The calculation supports a conditional use of model-based precision and a
clear reason to test the observation and state assumptions before relying
on it. An outcome audit, better follow-up or a richer measured history can
change those assumptions' credibility. More bootstrap draws alone cannot.
These results remain a generated research comparison with known reward
and initial laws; a bank application needs its own economic specification,
sampling argument and measurement evidence.

